\documentclass{standalone}
\usepackage{polymer-paper-preamble}
\begin{document}
\begin{tikzpicture}

%% --- Layout constants ---
%% Five panels arranged left-to-right in a 16 cm × 3.5 cm canvas.
%% Panel x-ranges: P1 0.3-3.0, P2 3.5-5.8, P3 6.3-9.0, P4 9.5-13.0, P5 13.5-15.8.
%% Panel title baseline y=3.55. Panel content baseline y=0.5..3.0.

%% =====================================================================
%% Panel 1 — Experimental trace (log-log power-law)
%% =====================================================================
\node[font=\sffamily\fontsize{7}{8.4}\selectfont, align=center, text=cGray!90!black]
  at (1.65, 3.55) {Experimental trace};

\LogLogPlot{0.45,0.6,3.30,3.0}%
  {0.55,1.10,1.65,2.20,2.75,3.20}%
  {0.55/10^{-3},1.10/10^{-1},1.65/10^{1},2.20/10^{2},2.75/10^{3},3.20/10^{4}}%
  {0.85,1.45,2.05,2.65}%
  {0.85/10^{-3},1.45/10^{-1},2.05/10^{1},2.65/10^{3}}%
  {$\log\,t$}{$\log\,I$}

% Power-law line (slope -n on log-log)
\draw[cBlue, line width=0.85pt] (0.55, 2.85) -- (3.20, 0.70);

% Slope triangle annotation
\draw[cBlue, dashed, line width=0.4pt] (1.20, 2.30) -- (1.80, 2.30);
\draw[cBlue, dashed, line width=0.4pt] (1.80, 2.30) -- (1.80, 1.80);

\node[font=\sffamily\fontsize{6}{7.2}\selectfont, text=cBlue]
  at (2.55, 2.55) {$I(t) \propto t^{-n}$};
\node[font=\sffamily\fontsize{5.5}{6.6}\selectfont, text=cBlue]
  at (1.30, 2.05) {slope $n$};

%% =====================================================================
%% Inter-panel arrow 1 → 2
%% =====================================================================
\draw[cGray, line width=0.55pt, -{Stealth[length=3.5pt,width=2.6pt]}]
  (3.40, 1.65) -- (3.65, 1.65);

%% =====================================================================
%% Panel 2 — Mathematical interpretation
%% =====================================================================
\node[font=\sffamily\fontsize{7}{8.4}\selectfont, align=center, text=cGray!90!black]
  at (4.65, 3.55) {Mathematical\\interpretation};

% Three extracted variables
\node[font=\sffamily\fontsize{8}{9.6}\selectfont, text=cBlue]
  at (4.00, 2.55) {$n$};
\node[font=\sffamily\fontsize{8}{9.6}\selectfont, text=cBlue]
  at (4.00, 1.65) {$\tau_d$};
\node[font=\sffamily\fontsize{8}{9.6}\selectfont, text=cBlue!45!cRed]
  at (4.00, 0.75) {$g(E_t)$};

% Variable arrows leading to chemical/physical origin panel
\draw[cGray, line width=0.45pt, -{Stealth[length=3pt,width=2.2pt]}]
  (4.25, 2.55) -- (5.75, 2.55);
\draw[cGray, line width=0.45pt, -{Stealth[length=3pt,width=2.2pt]}]
  (4.25, 1.65) -- (5.75, 1.65);
\draw[cGray, line width=0.45pt, -{Stealth[length=3pt,width=2.2pt]}]
  (4.25, 0.75) -- (5.75, 0.75);

% Debye exp(-t/τ) inset, attached to τ_d row
\begin{scope}[shift={(4.55, 1.05)}, scale=0.55]
  \draw[cGray, line width=0.32pt, -{Stealth[length=2pt,width=1.5pt]}]
    (0,0) -- (1.45,0);
  \draw[cGray, line width=0.32pt, -{Stealth[length=2pt,width=1.5pt]}]
    (0,0) -- (0,1.10);
  \draw[cGray!85!black, line width=0.55pt, smooth, samples=40, domain=0.05:1.40]
    plot (\x, {0.95*exp(-1.65*\x)});
  \node[font=\sffamily\fontsize{4.5}{5.4}\selectfont, text=cGray!90!black, align=center]
    at (0.92, 0.92) {Debye\\exp$(-t/\tau)$};
  \node[font=\sffamily\fontsize{4.3}{5.2}\selectfont, text=cGray!90!black]
    at (0.55, -0.20) {$\log\,t$};
  \node[font=\sffamily\fontsize{4.3}{5.2}\selectfont, text=cGray!90!black, rotate=90]
    at (-0.20, 0.55) {$\log$ response};
  \draw[cGray, dashed, line width=0.28pt] (0.45,0) -- (0.45,0.42);
  \node[font=\sffamily\fontsize{4.3}{5.2}\selectfont, text=cGray!90!black]
    at (0.50, -0.10) {$\tau_d$};
\end{scope}

%% =====================================================================
%% Inter-panel arrow 2 → 3
%% =====================================================================
\draw[cGray, line width=0.55pt, -{Stealth[length=3.5pt,width=2.6pt]}]
  (5.90, 1.65) -- (6.25, 1.65);

%% =====================================================================
%% Panel 3 — Chemical / physical origin (polymer chains with S markers)
%% =====================================================================
\node[font=\sffamily\fontsize{7}{8.4}\selectfont, align=center, text=cGray!90!black]
  at (7.65, 3.55) {Chemical / physical origin};

% Top chain row — sparse isolated S atoms (physical origin baseline)
\WavyChain{6.40, 8.95, 2.85, zigzag, {6.65, 7.20, 7.80, 8.45}, {}}

% Middle chain row — S-rich segment cluster + isolated, dashed bracket on cluster
\WavyChain{6.40, 8.95, 1.85, zigzag, {6.65, 6.85, 7.10, 7.30, 8.30, 8.65}, {}}
% Dashed cluster bracket (chemical origin highlight)
\draw[cGray, dashed, line width=0.4pt, rounded corners=2pt]
  (6.55, 1.55) rectangle (7.45, 2.45);
\node[font=\sffamily\fontsize{5.5}{6.6}\selectfont, text=cAmber, align=center]
  at (7.00, 1.20) {S-rich segments};

% Bottom chain row — single S atoms (physical origin reference)
\WavyChain{6.40, 8.95, 0.85, zigzag, {6.70, 7.65, 8.45}, {}}

% Callout arrows: chemical origin (left) and physical origin (right) of middle row
% Both labels constrained to Panel 3 x-range (6.40..8.95) to avoid Panel 2 overlap.
\draw[cGray, line width=0.35pt, -{Stealth[length=2.5pt,width=1.8pt]}]
  (6.85, 2.45) -- (6.55, 3.20);
\node[font=\sffamily\fontsize{5.5}{6.6}\selectfont, text=cGray!90!black, anchor=south east]
  at (6.95, 3.20) {chemical origin};
\draw[cGray, line width=0.35pt, -{Stealth[length=2.5pt,width=1.8pt]}]
  (7.55, 2.45) -- (8.10, 3.20);
\node[font=\sffamily\fontsize{5.5}{6.6}\selectfont, text=cGray!90!black, anchor=south west]
  at (7.95, 3.20) {physical origin};

%% =====================================================================
%% Inter-panel arrow 3 → 4
%% =====================================================================
\draw[cGray, line width=0.55pt, -{Stealth[length=3.5pt,width=2.6pt]}]
  (9.05, 1.65) -- (9.45, 1.65);

%% =====================================================================
%% Panel 4 — Trap-depth inference (band diagram with full Conduction Band /
%% Valence Band labels per briefing §2 vocabulary requirement)
%% =====================================================================
\node[font=\sffamily\fontsize{7}{8.4}\selectfont, align=center, text=cGray!90!black]
  at (11.30, 3.55) {Trap-depth inference};

% Teal rounded frame around band region
\draw[bandFrame] (9.55, 0.40) rectangle (13.05, 3.30);

% Band boxes — explicit Conduction Band / Valence Band labels (no abbreviations)
\node[bandbox, minimum width=2.50cm, minimum height=0.45cm] at (11.30, 2.95) {};
\node[font=\sffamily\fontsize{6}{7.2}\selectfont, text=cGray!90!black]
  at (11.30, 2.95) {Conduction Band};

\node[bandbox, minimum width=2.50cm, minimum height=0.45cm] at (11.30, 0.75) {};
\node[font=\sffamily\fontsize{6}{7.2}\selectfont, text=cGray!90!black]
  at (11.30, 0.75) {Valence Band};

% Vertical Energy axis on left interior
\draw[bandAxis] (9.85, 1.05) -- (9.85, 2.65);
\node[rotate=90, text=cGray!90!black, font=\sffamily\fontsize{6}{7.2}\selectfont]
  at (9.65, 1.85) {Energy};

% Shallow trap levels (amber) — slightly below CB
\TrapLevel{10.45}{2.45}{trapShallow}
\TrapLevel{10.95}{2.45}{trapShallow}
\TrapLevel{11.65}{2.40}{trapShallow}
\node[font=\sffamily\fontsize{5.5}{6.6}\selectfont, text=cAmber!85!black]
  at (12.10, 2.42) {$\cdots$};
\node[font=\sffamily\fontsize{5.5}{6.6}\selectfont, text=cAmber, font=\itshape]
  at (12.78, 2.42) {shallow};

% Deep trap levels (purple) — slightly above VB
\TrapLevel{10.45}{1.30}{trapDeep}
\TrapLevel{10.95}{1.30}{trapDeep}
\TrapLevel{11.65}{1.25}{trapDeep}
\node[font=\sffamily\fontsize{5.5}{6.6}\selectfont, text=cBlue!45!cRed!85!black]
  at (12.10, 1.27) {$\cdots$};
\node[font=\sffamily\fontsize{5.5}{6.6}\selectfont, text=cBlue!45!cRed, font=\itshape]
  at (12.78, 1.27) {deep};

%% =====================================================================
%% Panel 4 → Panel 5: three teal curved arrows (CB → shallow lobe,
%% middle → midline, Valence Band → deep lobe). 1:1 mapping per §6 invariant.
%% =====================================================================
\draw[cTeal, line width=0.7pt, -{Stealth[length=3pt,width=2.2pt]}]
  (13.10, 2.95) to[bend left=18] (13.85, 2.55);
\draw[cTeal, line width=0.7pt, -{Stealth[length=3pt,width=2.2pt]}]
  (13.10, 1.85) to[bend left=4] (13.85, 1.85);
\draw[cTeal, line width=0.7pt, -{Stealth[length=3pt,width=2.2pt]}]
  (13.10, 0.75) to[bend right=18] (13.85, 1.15);

%% =====================================================================
%% Panel 5 — Trap distribution (ISPD)
%% Two side-oriented bell lobes; shallow lobe slightly smaller per §6 invariant.
%% =====================================================================
\node[font=\sffamily\fontsize{7}{8.4}\selectfont, align=center, text=cGray!90!black]
  at (14.85, 3.55) {Trap distribution\\(ISPD)};

% Shallow lobe (amber) — top, smaller width
\BellCurve{13.95, 2.05, 15.20, 2.95, cAmber, side}
\node[font=\sffamily\fontsize{5.5}{6.6}\selectfont, text=cAmber, font=\itshape]
  at (15.55, 2.50) {shallow};

% Deep lobe (purple) — bottom, slightly larger width
\BellCurve{13.95, 0.65, 15.40, 1.85, cBlue!45!cRed, side}
\node[font=\sffamily\fontsize{5.5}{6.6}\selectfont, text=cBlue!45!cRed, font=\itshape]
  at (15.55, 1.25) {deep};

% g(E_t) horizontal axis at bottom
\draw[cGray, line width=0.4pt, -{Stealth[length=2.5pt,width=1.8pt]}]
  (13.90, 0.55) -- (15.95, 0.55);
\node[font=\sffamily\fontsize{6}{7.2}\selectfont, text=cGray!90!black]
  at (15.85, 0.30) {$g(E_t)$};

% g(E_t) top label (matches reference image)
\node[font=\sffamily\fontsize{6.5}{7.8}\selectfont, text=cTeal!75!black]
  at (15.85, 3.10) {$g(E_t)$};

\end{tikzpicture}
\end{document}
